\section{Introduction}
\label{secIntroduction}

Ce document rassemble et structure les \textbf{conventions de nommage} rencontrées dans le
développement d'une application \textbf{Control Expert} pilotant une cellule robotisée
associant un automate \textbf{Schneider M340}, un robot \textbf{Stäubli Scara} via le serveur
client \textbf{uniVALplc}, un contrôleur \textbf{M221} et un pupitre opérateur \textbf{HMI}.

Les conventions présentées ici sont \textbf{issues} de cinq documents
pédagogiques source (\texttt{RIF\_GT\_TP1\_2024}, \texttt{RIF\_RBT\_TP2\_2024},
\texttt{RIF\_RBT\_TP3\_2024}, \texttt{RIF\_RBT\_TP4\_2024}, \texttt{RIF\_RBT\_TP5\_2024}) rédigés
par un autre enseignant pour une cellule similaire. Elles ne constituent pas des énoncés de TP :
seules les conventions de nommage elles-mêmes (préfixes, types, blocs fonctionnels, sections)
ont été retenues et réorganisées par catégorie, indépendamment du fil pédagogique des TP dont
elles sont issues.

\UPSTIremarque[Allègement de la convention]{%
La convention d'origine, inspirée de la \textbf{notation hongroise} (préfixe \texttt{m\_} suivi
d'une ou plusieurs lettres de type, préfixes \texttt{nix}/\texttt{nqx}, instances
\texttt{Inst<Bloc><Mode>}), a été \textbf{allégée} : le type d'une variable n'est plus encodé dans
son nom, l'environnement de développement l'affichant à la complétion et au survol. Ce choix
s'aligne sur les pratiques actuelles (PLCopen, Codesys, TwinCAT) : des noms métier courts en
camelCase. Les types (\texttt{T<Nom>}/\texttt{T\_<Nom>}), les objets métier
(\texttt{<rôle><Équipement>}) et les sections/transitions du GEMMA sont inchangés. La table de
correspondance ancien/nouveau nom figure dans \texttt{MIGRATION\_NOTATION.md}.}

\medskip
Chaque tableau précise, pour une catégorie donnée, la convention, sa signification et un exemple
conforme à la convention allégée.
